%% E-LIS-AST-MET-1012.tex
%%
%% Front page for Critical algorithms.
%%
%% 2023-07-27: Version for FDR.
%%
\newcommand{\doctitle}{%
  \parbox[t]{\textwidth}{Software Modules,\\Simulated Data, and\\Documentation Package}}
\newcommand{\docnumber}	{E-LIS-AST-MET-1012}     % Document number
\newcommand{\issuenumber}{1-0}                % Version
\newcommand{\issuedate}{2023-07-27}              % Date
\newcommand{\workpackage}{8.2}      % Workpackage

\newcommand{\authorname}{%                       % Authors
  T.~Marquart,\\
  G.~Otten,\\
  N.~Sabha,\\
  M.~Baláž,\\
  O.~Czoske,\\
  H.~Buddelmeijer,\\
  K.~Leschinski,\\
}
\newcommand{\authorsigndate}  {2023-06-31} % date of signature of author

\newcommand{\reviewername}     {%
}       % checked by name
\newcommand{\reviewersigndate} {} % date of signature of checker

\newcommand{\approvername}    {%
}       % name of the approver
\newcommand{\approvalsigndate}{} % date of signature of approver

\newcommand{\releasername}{%
}
\newcommand{\releasesigndate}{}  % date of release

% Title displayed in the header
\newcommand{\shorttitle}{Data Reduction Library Validation and Test Plan}


%%%%%%%%%%%%%%%%%%%%%%%%%%%%%%%%%%%%%%%%%%%%%%%%%%%%%%%%%%%%%%%%%%%%%%%%%%%%%
\documentclass[11pt,oneside,a4paper]{article}

%% Layout  --- TeX commands
\raggedbottom
\topmargin=-9mm
\headsep=0.5cm
\textheight=246mm
\textwidth=164mm
\hoffset=0cm
\oddsidemargin=-2mm
\parindent=0mm
\parskip=0.3em

\renewcommand{\floatpagefraction}{0.98} % separate float page does not work
\renewcommand{\bottomfraction}{1.0}

%% Fonts
\usepackage{mathptmx}
\usepackage[scaled=0.92]{helvet}

\usepackage[pdftex]{graphicx}
\graphicspath{{./figures/}}
\renewcommand{\bottomfraction}{1.0}

\usepackage{fancyhdr}
\usepackage{titlesec}
\usepackage{booktabs}
\usepackage{tabularx}
\usepackage{array}
\usepackage{longtable}
\usepackage{numprint}
\usepackage{placeins}
\usepackage{enumitem}
\usepackage{acronym}
\usepackage{filecontents}

\usepackage{xcolor}
\usepackage{colortbl}
\definecolor{listingbg}{gray}{0.95}
\definecolor{darkgreen}{rgb}{0.0, 0.7, 0.0}
\definecolor{darkblue} {rgb}{0.0, 0.0, 0.7}
\definecolor{darkred}  {rgb}{0.7, 0.0, 0.0}
\definecolor{darkorange}{rgb}{1.0, 0.49, 0.0}

\usepackage[many]{tcolorbox}

% does not work with my version of biblatex:
% \usepackage[defernumbers=true,backend=biber,hyperref=true,url=false]{biblatex}
\usepackage[%
   defernumbers=true,
   backend=biber,
   hyperref=true,
   url=false,
   sorting=none]{biblatex}

\usepackage[%
% pagebackref=true,   % does not work with biblatex
   pdfpagelabels=true,
   plainpages=false,
   colorlinks=false]{hyperref}

% does not work with my version:
%\usepackage[xindy,nopostdot,nogroupskip]{glossaries}
\usepackage[xindy]{glossaries}
%\makeglossaries

\usepackage{lastpage}
\usepackage[%
   font=small,
   format=plain,
   indention=1.5em,
   labelfont={small,bf},
   up,
   justification=justified,
   singlelinecheck=true]{caption}

%% for sidewaysfigure and sidewaystable
\usepackage{rotating}

%% put recipe names, QC parameters, etc., in \lstinline{}
\usepackage{listings}
\lstloadlanguages{csh, c}

\usepackage{textcomp}

\usepackage{subfloat}

%% Write recipe names like this: \REC{metis_do_stuff}
\lstdefinestyle{RECstyle}{%
  basicstyle=\ttfamily\color{darkgreen}%
  \lst@ifdisplaystyle\scriptsize\fi}
\newcommand{\REC}[1]{\lstinline[style=RECstyle]!#1!}

%% Write requirements like this: \REQ{METIS-xxxx}
\newcommand{\REQ}[1]{\href{https://polarion.astron.nl/polarion/\#/project/METIS/workitem?id=#1}{\textcolor{brown}{#1}}}

%% Write code snippets like this: \CODE{SOMETHING==ELSE}
\newcommand{\CODE}[1]{\lstinline[]!#1!}

%% Use environment recipedef to define a new recipe.
\newenvironment{recipedef}
{% begin code
  \begin{tcolorbox}[breakable,enhanced jigsaw]%
    \begin{longtable}{p{0.24\hsize}p{0.74\hsize}}}
      {% end code
    \end{longtable}%
  \end{tcolorbox}}

%%%%% This is a definition that has the flow diagram within the box:
%% \newcommand{\theimage}{}  % dummy definition
%% \newenvironment{recipedef}[1]
%% {% begin
%% code
%% \renewcommand{\theimage}{#1}
%% \begin{tcolorbox}%
%%   \begin{tabular}{p{0.24\hsize}p{0.74\hsize}}
%%   }%
%%     {% end code
%%   \end{tabular}%
%%   \begin{center}%
%%     \includegraphics[width=0.6\textwidth]{\theimage}%
%%   \end{center}%
%% \end{tcolorbox}}


%% Proper typography for micrometer
\usepackage{upgreek}
\newcommand{\micron}{\upmu\mathrm{m}}

%% Ion
\DeclareRobustCommand{\ion}[2]{%
  \textup{#1{\mdseries\textsc{#2}}}%
}

%%%% Degrees, arcmin, arcsec
\newcommand{\degr}{\hbox{$^\circ$}}

\newcommand{\arcmin}{\hbox{$^\prime$}}

\newcommand{\arcsec}{\hbox{$^{\prime\prime}$}}

%% Bright red for things that still need to be looked at. None of
%% these should appear in the release version
\newcommand{\TODO}[1]{\textcolor{red}{\bfseries TODO: #1}}

\newcommand{\intentblankpage}{}

%% paragraph column, centred text
\newcolumntype{P}[1]{>{\centering\arraybackslash}p{#1}}

%=== Header / Footer definition ===========================
\newcommand{\headerformat}{%
  \lhead{\small\sffamily
    \begin{tabularx}{\textwidth}{Xlll}
      \shorttitle & \docnumber & \issuenumber & \issuedate\\[0.7ex]
    \end{tabularx}}
  \rhead{}
  \cfoot{\small\bfseries\sffamily
    Page {\textbf{\thepage}}  of {\textbf{\pageref{LastPage}}} }
  \rfoot{\raisebox{-0.3\height}{\includegraphics[width=25mm]{metis_noText}} }
}

\fancypagestyle{plain}{\fancyhf{} \headerformat}
\headerformat
\definecolor{brn}{RGB}{99,36,35}
\definecolor{rd1}{RGB}{229,184,183}
\definecolor{rd2}{RGB}{242,219,219}

\renewcommand\headrule{%
  \begingroup
  \color{brn}
  \hrule height 0.5pt width\headwidth
  \endgroup
}

\renewcommand{\contentsname}{Table of Contents}

%% ---- Format of section titles --------------------------------
% Combining sffamily and scshape works only for some fonts
\titleformat{\section}{\LARGE\sffamily\scshape}{\thesection}{1em}{\textsc{}}
\titleformat{\subsection}{\Large\sffamily}{\thesubsection}{1em}{}
\titleformat{\subsubsection}{\large\sffamily}{\thesubsubsection}{1em}{}

\setlength\heavyrulewidth{1.5pt}


% === PDF Definition ==========================================================
\hypersetup{%
  setpagesize,
  bookmarksnumbered=true,
  pdfborder= 0 0 0,
  pdftitle={\shorttitle},
  pdfauthor={\authorname}}


% === biblatex Definitions ====================================================
\begin{filecontents}{references.bib}
@Manual{DRLD,
  title =        {METIS Data Reduction Library Design},
  version =      {1-0},
  number =       {E-REP-AST-MET-1006},
  keywords =     {applicable},
  note =         {2023-09-18}
  }
\end{filecontents}

\bibstyle{apj}
\addbibresource{references.bib}
\newbibmacro{string+doiurlisbn}[1]{\iffieldundef{url}{#1}{\href{\thefield{url}}{#1}}}
\DeclareFieldFormat{title}{\usebibmacro{string+doiurlisbn}{\mkbibemph{#1}}}


% === glossaries definitions  =================================================
% \input{../common/metis_glossary.tex}     % the common glossary file
\renewcommand{\glossarysection}[2][]{}   % Avoid the heading 'Glossary'
\setlength{\LTleft}{0pt}                 % make all longtables aligned left
\setlength{\glsdescwidth}{0.8\textwidth}


%%%%%%%%%%%%%%%%%%%%%%%%%%%%%%%%%%%%%%%%%%%%%%%%%%%%%%%%%%%%%%%%%%%%%%%%%%%%%
%
% Start of document
%
%%%%%%%%%%%%%%%%%%%%%%%%%%%%%%%%%%%%%%%%%%%%%%%%%%%%%%%%%%%%%%%%%%%%%%%%%%%%%
%
\begin{document}

% listings style
\lstset{basicstyle=\ttfamily\lst@ifdisplaystyle\scriptsize\fi,
  columns=flexible,
  frame=single,
  backgroundcolor=\color{listingbg},
  captionpos=b,
  showspaces=false}


%%%%% titlepage
\thispagestyle{empty}

% Logo
\vspace*{0cm}
\includegraphics[width=6.69cm]{metis_logo}

\vspace*{\fill}

% Document title
{\color{brn}\rule[1.9ex]{\textwidth}{1.5pt}}
\scalebox{1.44}{\Huge\textsf{\doctitle}}\\
{\color{brn}\rule{\textwidth}{1.5pt}}\\ [0.5ex]

% Document info
{\Large\textsf{\docnumber}}  \\ [1ex]
{\Large\textsf{Issue \issuenumber}} \\ [1ex]
{\Large\textsf{\issuedate}}  \\ [1ex]
{\Large\textsf{Work package: \workpackage}}  \\[1ex]

\vspace*{\fill}

% Signature table
\begin{center}
  \renewcommand{\arraystretch}{0.75}
%  \begin{tabularx}{\textwidth}{XXXX}
  \begin{tabular}{p{0.25\textwidth}p{0.25\textwidth}p{0.25\textwidth}p{0.25\textwidth}}
    \arrayrulecolor{brn}
    \toprule
    & \multicolumn{3}{l}{\scriptsize\textsf{Signature and Approval}} \\
    \midrule
    & {\scriptsize\textsf Name}
    & {\scriptsize\textsf Date}
    & {\scriptsize\textsf Signature} \\
    \cline{2-4}
    \\
    \textsf{Prepared} & \parbox[c]{\hsize}{\raggedright \authorname} & \authorsigndate & \\
    \\
    \midrule
    \\
    \textsf{Reviewed} & \parbox[c]{\hsize}{\raggedright \reviewername} & \reviewersigndate & \\
    \\
    \midrule
    \\
    \textsf{Approved} & \parbox[c]{\hsize}{\raggedright \approvername} & \approvalsigndate & \\
    \\
    \midrule
    \\
    \textsf{Released} & \parbox[c]{\hsize}{\raggedright \releasername} & \releasesigndate & \\
    \\
    \bottomrule
     %    \end{tabularx}
  \end{tabular}
\end{center}

%%% Local Variables:
%%% TeX-master: "METIS_DRLVT"
%%% End:


% === MAIN TEXT================================================================

\clearpage
\pagestyle{fancy}

% === Revision History =========================================================
\section*{Revision History}
\renewcommand{\arraystretch}{1.2}
\arrayrulecolor{brn}
\begin{tabularx}{\textwidth}{|l|l|l|X|}
  \hline
  \rowcolor{rd1}
  \textbf{Issue} & \textbf{Date} & \textbf{Owner} & \textbf{Changes} \\
  \hline
  1-0 & 2023-07-27 & K. Leschinski & FDR release\\
  \hline
\end{tabularx}


%=== Indexes ==================================================================
%\newpage
%\tableofcontents
%\clearpage
%\listoffigures
%\clearpage
%\listoftables

%\clearpage
%\phantom{a}
%\vfill
%\begin{center}
%  This page intentionally left blank
%\end{center}
%\vfill
%\clearpage

%%%%%%%%%%%%%%%%%%%%%%%%%%%%%%%%%%%%%%%%%%%%%%%%%%%%%%%%%%%%%%%%%%%%%%%%%%%%%


% INTRODUCTION
\section{Introduction}
\label{sec:intro}

\subsection{Scope}
The critical algorithms as described in section 10 of the METIS Data Reduction Library Design \cite{DRLD} are stored in a git repository at \url{https://github.com/AstarVienna/CriticalAlgorithms}.

Only four of the critical algorithms require prototyping, as described in \cite{DRLD}.
These are
\begin{enumerate}
    \item[5a.] Long Slit Spectroscopy (LSS) Wavelength calibration and distortion correction.\\
        See directory \CODE{5a-LSS-wavelength-calibration-PyReduce-ELT}.
    \item[5b.] Integral Field Unit (IFU) Wavelength calibration and distortion correction.\\
        See directory \CODE{5b-IFU-distortion}.
    \item[9.] IFU image and cube reconstruction.\\
        See directory \CODE{9-IFU-image-reconstruction}.
    \item[11.] Angular Differential Imaging algorithm.\\
        See directory \CODE{11-ADI-algorithm}.
\end{enumerate}

\subsection{Applicable documents}
\begin{refcontext}[labelprefix=AD]
  \printbibliography[keyword=applicable, heading=none]
\end{refcontext}

%\subsection{Reference documents}
%\begin{refcontext}[labelprefix=RD]
%  \printbibliography[keyword=reference, heading=none]
%\end{refcontext}


%\clearpage

%\subsection{Acronyms}
%\label{ssec:acronyms}
%\input{acronyms.tex}

\clearpage



\end{document}

%%
%% THE END
%%
%%%%%%%%%%%%%%%%%%%%%%%%%%%%%%%%%%%%%%%%%%%%%%%%%%%%%%%%%%%%%%%%%%%%%%%%%%%%%
